\begin{textsample}{NEG Dim 3 – mt – Score -3.00 – mt/mt\_em\_48.txt}  \label{ex:f3_neg_203}
8.5 Estrutura da Área de Linguagens e suas Tecnologias para a Etapa Ensino Médio O Documento de Referência Curricular de Mato Grosso apresenta a área de Linguagens dividida em : \textbf{Competências} Específicas da Área e suas Habilidades.
Apenas o componente Língua Portuguesa traz suas \textbf{Competências} e Habilidades descritas.
Apresentam -se sete ( 7 ) \textbf{competências} específicas de linguagens e suas tecnologias.
Para cada \textbf{competência} são indicadas habilidades a serem desenvolvidas, ampliadas e aprofundadas ao longo do processo de ensino e aprendizagem do Ensino Médio ( BRASIL, 2018 ).
As \textbf{Competências} Específicas da área são as que devem ser desenvolvidas em todos os Componentes Curriculares ao longo do Ensino Médio, tanto na base comum como nos itinerários formativos.
São essas \textbf{Competências} que promovem a integração dos componentes da Área e a integralização da formação.
Essa busca por estabelecer uma integralidade da Área possibilita, por meio das práticas sociais, o desenvolvimento de \textbf{competências} e habilidades que ampliam as possibilidades de participação efetiva dos estudantes no mundo.
O DRC-MT – Ensino Médio apresenta também um Quadro comparativo em que são apresentadas as progressões do Ensino Fundamental ao Ensino Médio.
Importante notar quais os pontos de aprofundamento são acrescidos nas \textbf{competências} dessa etapa.
As \textbf{competências} específicas de Linguagens suas Tecnologias para o Ensino Médio reforçam e ampliam as já descritas para o Ensino Fundamental, como se pode observar no quadro abaixo :

% matched lemmas: competência
\end{textsample}
